% ==========================================================================
%  Template for a new research paper in this repository.
%  1. Copy this folder to research/NN-your-title/paper/
%  2. Fill in the title block and write sections/*.tex
%  3. Let the analysis code write paper/generated/values-*.tex
%  See research/STYLE_GUIDE.md for the full conventions.
% ==========================================================================
\documentclass[10pt,twocolumn,letterpaper]{article}
\usepackage{ewf-paper}

% Numbers written by the analysis code (uncomment when they exist)
% \input{generated/values-01-example}

% Path to the figure data, relative to this folder
\newcommand{\datadir}{../data/processed}

\begin{document}

\ewfmaketitle
  {Title of the Paper}
  {A Subtitle That States the Scope}
  {Author Name\footnotemark}
  {}% optional affiliation
  {Month Day, Year}
  {Abstract of 200 to 300 words. State the question, the method and the results with numbers.}
  {keyword one; keyword two; keyword three}
\footnotetext{This paper is a research compendium: figures read their numbers from \texttt{data/processed}, numbers in the text are written by \texttt{code/}, and external facts are listed with sources in \texttt{data/literature}.}

\section{Thesis}
State the problem and the claim of the paper in the first paragraphs.

\section{Example figure}
\begin{figure}[t]
\centering
% \readcsv{\datadir/fig01_example.csv}\exampledata
\begin{tikzpicture}
\begin{axis}[paper, title={Example}, xlabel={$x$}, ylabel={$y$}]
\addplot+[no marks,domain=0:10] {x};
\addplot+[no marks,domain=0:10] {0.5*x+1};
\addplot[reference,domain=0:10,samples=2] {5};
\end{axis}
\end{tikzpicture}
\caption{\textbf{The first sentence states the finding.} Then say what is shown and where the data are.}
\label{fig:example}
\end{figure}

\begin{proposition}[Name]
\label{prop:example}
Statement.
\end{proposition}
\begin{proof}
Argument.
\end{proof}

{\small
\bibliographystyle{unsrtnat}
% \bibliography{references}
}
\end{document}
